\documentclass[11pt]{article}
\usepackage[margin=25mm]{geometry}
\usepackage{fontspec}
\setmainfont{Latin Modern Roman}
\usepackage{amsmath,amssymb,mathtools}
\usepackage{unicode-math}
\setmathfont{Latin Modern Math}
\usepackage{xurl,hyperref}
\hypersetup{colorlinks=true,urlcolor=blue}
\setcounter{secnumdepth}{0}
\setlength{\parindent}{0pt}
\setlength{\parskip}{5pt}
\setlength{\emergencystretch}{3em}
\newcommand{\sref}[2]{\hyperlink{source:#1:#2}{[S#2]}}
\newcommand{\cataloguescope}[1]{\paragraph{Recorded target.}#1}
\newcommand{\R}{\mathbb R}
\newcommand{\C}{\mathbb C}
\newcommand{\N}{\mathbb N}
\newcommand{\Z}{\mathbb Z}
\newcommand{\Q}{\mathbb Q}
\newcommand{\F}{\mathbb F}
\newcommand{\eps}{\varepsilon}
\begin{document}
% UnsolvedMath ID 20002521; source catalogue code AIM-PHYSICS-0029
\setcounter{section}{20002520}
\section{Melrose's answer, compact Sobolev embedding, and escape of mass}\label{20002521}
{\small UnsolvedMath ID 20002521 · Mathematical Physics}
\subsection{Definitions and mathematical statement}

\paragraph{Shared notation.}$\mathbb N=\{1,2,\ldots\}$; the standard number fields have their usual meanings. All additional conventions are stated locally.


\paragraph{Definitions and notation.}For a Riemannian manifold $(M,g)$, $L^2$ norms use the metric volume and pointwise metric on forms. Reduced $L^2$ cohomology is $\ker d/\overline{\operatorname{im}d}$ for the specified closed $L^2$ complex. On a complete manifold its harmonic representatives satisfy $d\alpha=d^*\alpha=0$ and $\int_M|\alpha|^2<\infty$. For a triangulated universal cover, square-summable chains have coefficients with $\sum|a_i|^2<\infty$, together with the domain condition needed for their boundary operator. Write $D=d+d^*$. In the source’s nonparabolic-at-infinity setting, a compact core $K\subset\widetilde K$ gives the norm $N_{\widetilde K}(\alpha)=\|\alpha\|_{L^2(\widetilde K)}+\|D\alpha\|_{L^2(M\setminus\widetilde K)}$; $W$ is the completion of compactly supported smooth forms in that norm.

\paragraph{Catalogue wording retained for traceability.}
Define $L^2$ homology in a manner compatible with the square-summable-chain motivation; determine its dependence on the metric and when it is dual to $L^2$ cohomology. Determine the relation with group cohomology, for which the source refers to Lück. In the specified nonparabolic-at-infinity setting, describe $W$ alternatively; realize reduced $L^2$ cohomology as the cohomology of an appropriate complex or complexes; and put $W$ into a complex where possible. Determine when a Mayer--Vietoris sequence exists for reduced $L^2$ cohomology. The source records that nonparabolicity alone is insufficient, and asks whether additional closed-range control of $d$ on noncompact overlaps is sufficient.
\paragraph{Scope and source qualification.}The first sentence of the excerpt is Melrose’s answer to the prior Sobolev question; it is not substituted for this record’s actual L² questions. Closed range on overlaps remains a proposed sufficient condition rather than an established theorem. The catalogue title/status was paired with a machine-generated research attempt; that attempt is not peer-reviewed evidence of settlement. The statement covers the original source questions rather than only the title’s subcase.
\subsection{Short English statement}
Define metric-sensitive L² homology, clarify the Fredholm completion and find conditions for reduced L² excision and Mayer–Vietoris sequences.
\subsection{Research attempt}
\paragraph{Model, date and investigated scope.}
OpenAI GPT-6 Astra (configured model \texttt{gpt-6-astra}), ultra reasoning; generated 11 October 2026. Exhibited dense nonclosed range and proved reduced-cohomology invariance under uniformly equivalent norms with fixed domains. The requested L2 homology, W-complex, duality and Mayer-Vietoris theory remain unresolved.

\paragraph{Sources and literature check.}
Read the supplied complete definition, catalogue wording, scope qualifications and source list for this canonical ID. The original-source citations are inherited as supplied references; no fresh primary-source or novelty search was performed for this attempt. No current open-status claim is inferred from this calculation.

\paragraph{Why closure cannot be omitted.}On \(\ell^2(\mathbb N)\), define the bounded operator \(T(x_n)=(x_n/n)\). It is injective and its range contains all finitely supported sequences, hence is dense. It is not onto: \(y_n=1/n\) lies in \(\ell^2\), but its formal preimage is the constant sequence, which does not. Thus the two-term complex \(0\to\ell^2\xrightarrow{T}\ell^2\to0\) has nonzero ordinary cokernel but zero reduced cokernel, since the latter divides by the closure of the image. Approximate exactness and exactness differ concretely.
\paragraph{A metric-comparison statement.}Suppose two Hilbert norms on every space of a fixed differential complex are uniformly equivalent, with the same closed-operator domains. Then they induce the same topology, the same kernels, the same range closures and therefore naturally isomorphic reduced cohomology spaces. Indeed a sequence converges in one norm exactly when it converges in the other. This is a precise sufficient metric-independence condition, not invariance under arbitrary geometric changes at infinity; unbounded metric distortions need not preserve the underlying \(L^2\) spaces or domains.
\paragraph{First gap and source correction.}For a Mayer--Vietoris construction one needs bounded restriction/extension maps on the actual closed domains and control of ranges when passing to the reduced quotient. The example explains why a purely algebraic gluing argument can fail. It does not prove that closed range on noncompact overlaps is sufficient in the source's setting. The square-summable homology definition, duality with cohomology, relation to group cohomology, alternative characterization of the Carron completion \(W\), and embedding \(W\) into a complex remain unresolved. The preceding Sobolev remark in the source is not substituted for these L2 questions.

\paragraph{Verification and adversarial audit.}
Checked dense range, explicit missing preimage, ordinary versus reduced cokernel, fixed-domain requirement and restored L2 rather than prior Sobolev scope. This was an LLM adversarial audit, not independent human verification; no theorem prover was used. Prompts A then B applied to this ID: Checked dense range, explicit missing preimage, ordinary versus reduced cokernel, fixed-domain requirement and restored L2 rather than prior Sobolev scope. Scope remains limited: Construct bounded gluing maps on the actual L2 domains and prove the range control required for reduced exactness.

Coordinator checked the dense nonclosed diagonal-operator range and invariance of reduced cohomology under the stated equivalent-domain norms. This additional check was model-assisted, not independent human verification.

\paragraph{Tools, process and contributions.}
Python 3 and local shell used to read supplied files and serialize this manuscript; mathematics derived symbolically in the text, with no numerical experiment or formal proof assistant. Contribution prompt A was followed by prompt B. GPT-6 Astra with ultra reasoning drafted the research and performed its stated self-audit; any additional coordinator contribution is identified in the verification and audit. A GPT-6-based Codex coordinator assembled the manuscript and checked repository integration; its exact variant and reasoning configuration were not exposed. No human mathematical contribution, independent human proof check or formal proof certification is supplied. Complete derivations and the consulted source references are included here; private reasoning transcripts are not published.

\paragraph{Final status and first remaining gap.}
\textbf{UNRESOLVED}. Construct bounded gluing maps on the actual L2 domains and prove the range control required for reduced exactness.

\section{Completion Estimate}
\noindent\textbf{COMPLETION: 0\%}
The calculation is an elementary benchmark or conditional reduction; it establishes no claimed new portion of the complete source target.

\subsection{Sources}
\begin{itemize}
\item[\hypertarget{source:20002521:1}{S1}] ULAM.AI and the UnsolvedMath Contributors, UnsolvedMath, record 20002521 (AIM-PHYSICS-0029). Catalogue curation under CC BY 4.0; definitions and English exposition added locally.
\url{https://huggingface.co/datasets/ulamai/UnsolvedMath}.
\item[\hypertarget{source:20002521:2}{S2}] American Institute of Mathematics, $L^2$ harmonic forms in geometry and string theory; notes by Anda Degeratu and Mark Haskins, March 2004, 17 March 2004, L2-cohomology and Carron-lecture questions.
\url{https://aimath.org/WWN/l2harmonic/l2harmonic.pdf}.
\end{itemize}
\label{end:20002521}
\end{document}
